\chapter{Complexity — Random Bits}
%%% generated by the blueprint pipeline from TCSlib.Complexity.Randomized.RandomBits
\section{Overview}

Declarations of the Lean module \texttt{TCSlib.Complexity.Randomized.RandomBits}.

\section{Declarations}

\begin{theorem}[Probability of a constant predicate]
\lean{Randomized.randProb_const}
\label{Randomized.randProb_const}
\leanok
\uses{Randomized.randProb, Randomized.randProb_congr, Randomized.randProb_false, Randomized.randProb_true}
\difficulty{4}
Let $n\in\mathbb{N}$ and let $p$ be a decidable proposition. The probability that a
uniformly random Boolean word of length $n$ satisfies the predicate that is constantly
$p$ is one if $p$ holds and zero otherwise.
\end{theorem}

\begin{theorem}[Positive Randomized Probability Has a Witness]
\lean{Randomized.exists_ofFn_of_randProb_pos}
\label{Randomized.exists_ofFn_of_randProb_pos}
\leanok
\uses{Randomized.randProb, Randomized.randProb_congr, Randomized.randProb_false}
\difficulty{3}
Let $m$ be a natural number and let $A$ be a decidable property of Boolean lists. If the
fraction of length-$m$ Boolean words $w$ for which $A(w)$ holds is positive, then some
function $r : 	ext{Fin}(m) \to \{\mathrm{false},\mathrm{true}\}$ has $A$ on its
associated length-$m$ list.
\end{theorem}

\begin{theorem}[Probability of a specified single bit]
\lean{Randomized.randProb_one_eq}
\label{Randomized.randProb_one_eq}
\leanok
\uses{Randomized.randProb}
\difficulty{3}
For either Boolean value $c$, a uniformly random Boolean string of length one equals the
one-bit string $[c]$ with probability $1/2$.
\end{theorem}

\begin{theorem}[One-bit recurrence for word probability]
\lean{Randomized.randProb_succ}
\label{Randomized.randProb_succ}
\leanok
\uses{Randomized.randProb, Randomized.randProb_congr, Randomized.randProb_one_eq, Randomized.randProb_or_disjoint, Randomized.randProb_split}
\difficulty{5}
Let $n\in\mathbb N$ and let $P$ be a predicate on Boolean words. The probability that a
uniformly random Boolean word of length $n+1$ satisfies $P$ is the average of the
probabilities that a uniformly random length-$n$ word, prefixed by true or by false,
satisfies $P$.
\end{theorem}

\begin{theorem}[Probability after dropping a prefix]
\lean{Randomized.randProb_drop}
\label{Randomized.randProb_drop}
\leanok
\uses{Randomized.randProb, Randomized.randProb_congr, Randomized.randProb_split, Randomized.randProb_true}
\difficulty{3}
Let $S,n\in\mathbb{N}$ and let $P$ be a predicate on Boolean words. The probability that
the suffix after the first $S$ bits of a uniformly random Boolean word of length $S+n$
satisfies $P$ equals the probability that a uniformly random Boolean word of length $n$
satisfies $P$.
\end{theorem}

\begin{theorem}[Probability of a predicate on a word segment]
\lean{Randomized.randProb_drop_take}
\label{Randomized.randProb_drop_take}
\leanok
\uses{Randomized.randProb, Randomized.randProb_congr, Randomized.randProb_split, Randomized.randProb_take, Randomized.randProb_true}
\difficulty{4}
Let $s,T,m\in\mathbb{N}$ satisfy $s+T\le m$, and let $B$ be a predicate on Boolean
words. The probability that the length-$T$ segment obtained by dropping the first $s$
bits of a uniformly random Boolean word of length $m$ satisfies $B$ equals the
probability that a uniformly random Boolean word of length $T$ satisfies $B$.
\end{theorem}

\begin{theorem}[Probability of an Appended Event]
\lean{Randomized.randProb_append}
\label{Randomized.randProb_append}
\leanok
\uses{Randomized.randProb}
\difficulty{4}
For natural numbers $b,m$, let $A$ be a property of finite Boolean lists. For a property
$D$ of Boolean lists, let $P_k(D)$ denote the fraction of Boolean lists of length $k$
satisfying $D$. Then $P_{b+m}(A)=2^{-b}\sum_{u\in\{0,1\}^b}P_m(l\mapsto
A(u\mathbin{+	hinspace+}l))$, where $u\mathbin{+	hinspace+}l$ denotes concatenation.
\end{theorem}

\begin{theorem}[A String Avoiding Every Bad Event]
\lean{Randomized.exists_forall_not_of_randProb_le}
\label{Randomized.exists_forall_not_of_randProb_le}
\leanok
\uses{Randomized.card_univ_bitstrings, Randomized.randProb}
\difficulty{5}
Let $m,n$ be nonnegative integers, and let $\mathrm{Bad}(v,w)$ be a predicate on an
$n$-bit string $v$ and a bit string $w$. Suppose that for every $v$, the fraction of the
$2^m$ bit strings $w$ of length $m$ for which $\mathrm{Bad}(v,w)$ holds is at most
$\delta$, and that $2^n\delta<1$. Then there is an $m$-bit string $w$ for which
$\mathrm{Bad}(v,w)$ fails for every $n$-bit string $v$.
\end{theorem}

\begin{theorem}[Counts of complementary block predicates]
\lean{Randomized.blockCount_add_not}
\label{Randomized.blockCount_add_not}
\leanok
\uses{Randomized.blockCount, Randomized.blockCount_succ}
\difficulty{4}
For natural numbers $q,K$, a predicate $B$ on finite bit strings, and a finite bit
string $l$, the number of the first $K$ consecutive length-$q$ blocks of $l$ satisfying
$B$ plus the number satisfying its negation is $K$.
\end{theorem}

\begin{definition}[Binary value of a list block]
\lean{Randomized.blockVal}
\label{Randomized.blockVal}
\leanok
\uses{Complexity.LogProg.bitsVal}
For natural numbers $b,i$ and a Boolean list $r$, take the block of $b$ entries of $r$
starting at position $i b$ and let its binary value be the natural number assigned to
that block.
\end{definition}

\begin{theorem}[Bound on a block value]
\lean{Randomized.blockVal_lt}
\label{Randomized.blockVal_lt}
\leanok
\uses{Complexity.LogProg.bitsVal_lt_two_pow, Randomized.blockVal}
\difficulty{1}
For natural numbers $b,i$ and a finite bit string $r$, the integer represented by the
length-at-most-$b$ block of $r$ beginning at position $ib$ is strictly less than $2^b$.
\end{theorem}

\begin{theorem}[Positional expansion of a bit string]
\lean{Randomized.bitsVal_ofFn}
\label{Randomized.bitsVal_ofFn}
\leanok
\uses{Complexity.LogProg.bitsVal}
\difficulty{4}
For every natural number $L$ and every bit string $f : \mathrm{Fin}(L) \to \{0,1\}$, the
value of the list of its bits, with the first entry least significant, is
$\sum_{j=0}^{L-1} f(j)2^j$.
\end{theorem}

\begin{definition}[Boolean vectors and bounded integers]
\lean{Randomized.bitsEquiv}
\label{Randomized.bitsEquiv}
\leanok
For every natural number $L$, there is a bijection between Boolean-valued functions on
$\{0,\ldots,L-1\}$ and the integers $0,1,\ldots,2^L-1$, given by binary encoding.
\end{definition}

\begin{theorem}[Binary value of a bit string]
\lean{Randomized.bitsEquiv_apply}
\label{Randomized.bitsEquiv_apply}
\leanok
\uses{Complexity.LogProg.bitsVal, Randomized.bitsEquiv, Randomized.bitsVal_ofFn}
\difficulty{3}
For every natural number $L$ and every bit string $f$ indexed by the positions
$0,1,\ldots,L-1$, the natural number represented by its image under the standard
bijection between length-$L$ bit strings and integers from $0$ to $2^L-1$ is the value
of the list of its bits, with the first list entry interpreted as the least significant
bit.
\end{theorem}

\begin{theorem}[Probability of a Property of Binary Values]
\lean{Randomized.randProb_bitsVal}
\label{Randomized.randProb_bitsVal}
\leanok
\uses{Complexity.LogProg.bitsVal, Complexity.LogProg.bitsVal_lt_two_pow, Randomized.bitsEquiv, Randomized.bitsEquiv_apply, Randomized.randProb}
\difficulty{4}
Let $L$ be a natural number and let $P$ be a decidable property of natural numbers.
Interpret a Boolean string $r$ of length $L$ as the natural number
$v(r)=\sum_{i=0}^{L-1}r_i2^i$. The fraction of length-$L$ Boolean strings for which
$P(v(r))$ holds equals $2^{-L}$ times the number of integers $v<2^L$ for which $P(v)$
holds.
\end{theorem}

\begin{definition}[Binary block encoding]
\lean{Randomized.blocksEquiv}
\label{Randomized.blocksEquiv}
\leanok
\uses{Randomized.bitsEquiv}
For natural numbers $n$ and $b$, there is a bijection between Boolean strings of length
$nb$ and length-$n$ sequences whose entries are integers from $0$ through $2^b-1$,
obtained by grouping the bits into $n$ consecutive blocks of length $b$ and interpreting
each block as a binary integer.
\end{definition}

\begin{theorem}[Value of a block under the block bijection]
\lean{Randomized.blocksEquiv_apply}
\label{Randomized.blocksEquiv_apply}
\leanok
\uses{Randomized.bitsEquiv_apply, Randomized.blockVal, Randomized.blocksEquiv}
\difficulty{4}
For natural numbers $n,b$, a bit string $r$ indexed by the $nb$ positions, and an index
$i$ among the first $n$ positions, the integer represented by the $i$th block under the
standard bijection between length-$b$ bit strings and integers from $0$ to $2^b-1$
equals the value of the length-at-most-$b$ substring of the list of bits of $r$
beginning at position $ib$.
\end{theorem}

\begin{theorem}[Probability of a block predicate]
\lean{Randomized.randProb_blocks}
\label{Randomized.randProb_blocks}
\leanok
\uses{Randomized.blockVal, Randomized.blocksEquiv, Randomized.blocksEquiv_apply, Randomized.randProb}
\difficulty{3}
Let $n,b\in\mathbb{N}$ and let $P$ be a predicate on functions from $\{0,\ldots,n-1\}$
to $\mathbb{N}$. The probability that a uniformly random Boolean word of length $nb$,
divided into $n$ consecutive blocks of length $b$, has block values satisfying $P$
equals the number of functions $f:\{0,\ldots,n-1\}\to\{0,\ldots,2^b-1\}$ satisfying $P$,
divided by $(2^b)^n$.
\end{theorem}

\begin{theorem}[Probability of a predicate on blocks]
\lean{Randomized.randProb_blocks_piFinset}
\label{Randomized.randProb_blocks_piFinset}
\leanok
\uses{Randomized.blockVal, Randomized.randProb, Randomized.randProb_blocks}
\difficulty{4}
Let $n,b\in\mathbb{N}$ and let $P$ be a predicate on functions from $\{0,\ldots,n-1\}$
to $\mathbb{N}$. The probability that a uniformly random Boolean word of length $nb$,
divided into $n$ consecutive blocks of length $b$, has block values satisfying $P$
equals the number of functions $f:\{0,\ldots,n-1\}\to\{0,\ldots,2^b-1\}$ satisfying $P$,
divided by $2^{nb}$.
\end{theorem}

\begin{theorem}[Probability of success in every block]
\lean{Randomized.randProb_forall_blocks}
\label{Randomized.randProb_forall_blocks}
\leanok
\uses{Complexity.LogProg.bitsVal, Randomized.blockVal, Randomized.randProb, Randomized.randProb_congr, Randomized.randProb_split, Randomized.randProb_true}
\difficulty{5}
Let $t,L$ be natural numbers and let $Q$ be a predicate on natural numbers. For a
uniformly random Boolean word of length $tL$, divide the word into $t$ consecutive
blocks of length $L$ and interpret each block as a binary integer. The probability that
every block's value satisfies $Q$ equals the probability that the value of one uniformly
random Boolean word of length $L$ satisfies $Q$, raised to the power $t$.
\end{theorem}

\begin{theorem}[Lower Bound for Appended-Event Probability]
\lean{Randomized.randProb_append_ge}
\label{Randomized.randProb_append_ge}
\leanok
\uses{Randomized.randProb, Randomized.randProb_nonneg}
\difficulty{5}
Let $m_1,m_2$ be natural numbers, let $A$ be a decidable property of Boolean lists, and
let $C(u,v)$ be a decidable property of pairs of Boolean lists. Write $P_k(D)$ for the
fraction of Boolean lists of length $k$ satisfying a property $D$. Suppose that for
every Boolean string $u$ of length $m_1$ satisfying $A$, the probability that an
independent uniform Boolean string $v$ of length $m_2$ satisfies $C(u,v)$ is at least
$\varepsilon$, where $\varepsilon$ is rational. Then the probability that a uniform
Boolean string $r$ of length $m_1+m_2$ satisfies $C(r_{\le m_1},r_{>m_1})$ is at least
$P_{m_1}(A)\varepsilon$.
\end{theorem}

\begin{theorem}[Upper Bound for Appended-Event Probability]
\lean{Randomized.randProb_append_le}
\label{Randomized.randProb_append_le}
\leanok
\uses{Randomized.randProb, Randomized.randProb_append_ge, Randomized.randProb_congr, Randomized.randProb_not, Randomized.randProb_or_disjoint, Randomized.randProb_true}
\difficulty{5}
Let $m_1,m_2$ be natural numbers, let $A$ be a decidable property of Boolean lists, and
let $C(u,v)$ be a decidable property of pairs of Boolean lists. Write $P_k(D)$ for the
fraction of Boolean lists of length $k$ satisfying a property $D$. Suppose that $C(u,v)$
implies $A(u)$ for all Boolean lists $u,v$, and that for every length-$m_1$ string $u$
satisfying $A$, the probability that a uniform length-$m_2$ string $v$ satisfies
$C(u,v)$ is at most the rational number $\varepsilon$. Then the probability that a
uniform Boolean string $r$ of length $m_1+m_2$ satisfies $C(r_{\le m_1},r_{>m_1})$ is at
most $P_{m_1}(A)\varepsilon$.
\end{theorem}
